\section{Conclusion}
\label{sec:limitations}

The covered-call representation unifies designated collection, history-dependent prediction
and own-state costs. It gives the independent value, optimal policy and exact action losses,
recovers established special cases, prices collection rights and costs, bounds the value of resumption, and
quantifies how posterior updates under sharing perturb those guarantees through
changes in local continuation values. A common-component envelope relates those
changes to predictive total-variation dispersion; its decay must be controlled
over all feasible future histories.

Exact optimality requires finite independent arm processes, arm-local availability and
costs, and readings that follow the planning law; general information sharing invalidates the
proof, and the forward-closure stopping condition is not established for shared surrogates.
The sharing bound requires a coherent joint path law, arm-local costs, true posterior
marginals and an envelope over every future feasible history; it is conservative in the exact
examples and certifies neither posterior accuracy nor numerical index error.
Weak posterior concentration alone does not imply that this envelope vanishes.
An unbounded Gaussian predictive law also requires an extension of the stated
bounded-law assumptions. Continuous
reading laws require integration or discretization whose policy error is not bounded here,
observation trees are exponential in depth, and the objective is the best observed
collectable score less cost, not latent performance under noisy readings. The studies use
two arms, few outcomes, short ladders and known laws; they measure neither learned-surrogate
accuracy nor robustness to a misspecified law and establish no empirical superiority for
hyperparameter optimization. Scaling, misspecification, continuous candidate optimization and
comparisons with published methods on their own surrogates remain open for the staged
index; \citet{vu2026coststop} report such comparisons for the one-step rule.
